\documentclass[tikz,border=10pt]{standalone}
\usepackage{ctex}
\usepackage{tikz}
\usetikzlibrary{arrows.meta,positioning}

\begin{document}
\begin{tikzpicture}[
  font=\small,
  >=Latex,
  step/.style={draw=black!70, fill=blue!8, rounded corners=2mm, minimum width=36mm, minimum height=10mm, align=center},
  io/.style={draw=black!70, fill=teal!10, rounded corners=2mm, minimum width=32mm, minimum height=10mm, align=center},
  line/.style={->, very thick}
]

\node[io] (eth) at (0,0) {以太网报文到达};
\node[step, right=18mm of eth] (push) {写入环形存储\\\(\texttt{push}(e_i)\)};
\node[io, below=12mm of push] (can) {CAN 异常触发};
\node[step, right=18mm of push] (query) {时间窗检索\\\(\texttt{items\_within\_timeframe}\)};
\node[step, right=18mm of query] (merge) {中央处理器融合\\构建告警载荷};
\node[io, right=18mm of merge] (out) {输出告警\\(含以太网上下文)};

\draw[line] (eth) -- (push);
\draw[line] (push) -- (query);
\draw[line] (can.east) |- (query.south);
\draw[line] (query) -- (merge);
\draw[line] (merge) -- (out);

\node[draw=black!35, rounded corners=2mm, fill=orange!12, align=left, inner sep=5pt] at (9.2,-2.0) {
触发逻辑：仅在 CAN 异常发生时启动检索\\
运行优势：写路径常量时延，读路径按需执行
};

\end{tikzpicture}
\end{document}
